% GENERATED FILE. Edit canonical lessons, then run npm run generate:books.
% canonical-source-hash: fnv1a64:13ace86fd24271f0
% canonical-lessons: FR-C439-cathedrale, FR-C439-basilique, FR-C439-clocher, FR-C439-beffroi, FR-C439-remparts, FR-R439-revision-432-439

\chapter{Clochers et remparts}
\label{ch:fr-clochers-et-remparts}
\hlchaptermodality{439}

\begin{chapteropening}
\textbf{By the end of this chapter:} \emph{I can talk about the old buildings of a town: the cathedral, the basilica, the bell tower, the belfry and the town walls.}

Complete the last lesson of chapter 439: I can talk about the old buildings of a town: the cathedral, the basilica, the bell tower, the belfry and the town walls.
\end{chapteropening}

\section[la cathédrale]{la cathédrale — the cathedral}

\label{lesson:FR-C439-cathedrale}



\begin{quote}
Before the new one: say the French for to wander, then the French for to pace up and down.
\end{quote}

\subsection*{You'll want to know: la cathédrale}
\textbf{la cathédrale} — \textquotedblleft{}the cathedral\textquotedblright{}.

\begin{quote}
La cathédrale est au centre de la vieille ville.
\end{quote}

\begin{quote}
\emph{The cathedral is in the middle of the old town.}
\end{quote}

The bishop's church.

\begin{cousinweb}
\emph{cathédrale} is a learned borrowing from Medieval Latin \emph{(ecclēsia) cathedrālis} 'church with a bishop's chair', from \emph{cathedra} 'chair', from Greek \emph{kathédra} (Lewis \& Short).
\end{cousinweb}

\begin{grammarlens}[title={What you've built}]
One more word for telling, explaining and describing at B1.
\end{grammarlens}

\subsection*{Your turn}
\begin{itemize}
\raggedright
  \item \emph{Say it:} \emph{la cathédrale}
  \item \emph{Say it:} \emph{la cathédrale}, once more
  \item \emph{Say it:} a sentence of your own with \emph{la cathédrale}
  \item \emph{Recall:} say the French for to wander, then the French for to pace up and down, then say \emph{la cathédrale} again
\end{itemize}

\begin{tcolorbox}[breakable,colback=teal!4,colframe=teal!35!black,title={Before you move on}]
What does la cathédrale mean? (\textquotedblleft{}The cathedral\textquotedblright{}.) Say it once more.
\end{tcolorbox}
\section[la basilique]{la basilique — the basilica}

\label{lesson:FR-C439-basilique}



\begin{quote}
Before the new one: say the French for to pace up and down, then the French for the cathedral.
\end{quote}

\subsection*{You'll want to know: la basilique}
\textbf{la basilique} — \textquotedblleft{}the basilica\textquotedblright{}.

\begin{quote}
On monte à la basilique par un long escalier.
\end{quote}

\begin{quote}
\emph{You go up to the basilica by a long flight of steps.}
\end{quote}

A large and important church.

\begin{cousinweb}
\emph{basilique} is a learned borrowing from Latin \emph{basilica}, from Greek \emph{basilik\'{\={e}} (stoá)} 'royal hall' (EtymDB).
\end{cousinweb}

\begin{grammarlens}[title={What you've built}]
One more word for telling, explaining and describing at B1.
\end{grammarlens}

\subsection*{Your turn}
\begin{itemize}
\raggedright
  \item \emph{Say it:} \emph{la basilique}
  \item \emph{Say it:} \emph{la basilique}, once more
  \item \emph{Say it:} a sentence of your own with \emph{la basilique}
  \item \emph{Recall:} say the French for to pace up and down, then the French for the cathedral, then say \emph{la basilique} again
\end{itemize}

\begin{tcolorbox}[breakable,colback=teal!4,colframe=teal!35!black,title={Before you move on}]
What does la basilique mean? (\textquotedblleft{}The basilica\textquotedblright{}.) Say it once more.
\end{tcolorbox}
\section[le clocher]{le clocher — the bell tower, the steeple}

\label{lesson:FR-C439-clocher}



\begin{quote}
Before the new one: say the French for the cathedral, then the French for the basilica.
\end{quote}

\subsection*{You'll want to know: le clocher}
\textbf{le clocher} — \textquotedblleft{}the bell tower, the steeple\textquotedblright{}.

\begin{quote}
On voit le clocher du village de très loin.
\end{quote}

\begin{quote}
\emph{The village steeple is visible from far away.}
\end{quote}

\emph{L'esprit de clocher} — small-town thinking.

\begin{cousinweb}
\emph{clocher} is built on \emph{cloche} 'bell', which goes back to Late Latin \emph{clocca}, usually traced to a Celtic word.
\end{cousinweb}

\begin{grammarlens}[title={What you've built}]
One more word for telling, explaining and describing at B1.
\end{grammarlens}

\subsection*{Your turn}
\begin{itemize}
\raggedright
  \item \emph{Say it:} \emph{le clocher}
  \item \emph{Say it:} \emph{le clocher}, once more
  \item \emph{Say it:} a sentence of your own with \emph{le clocher}
  \item \emph{Recall:} say the French for the cathedral, then the French for the basilica, then say \emph{le clocher} again
\end{itemize}

\begin{tcolorbox}[breakable,colback=teal!4,colframe=teal!35!black,title={Before you move on}]
What does le clocher mean? (\textquotedblleft{}The bell tower, the steeple\textquotedblright{}.) Say it once more.
\end{tcolorbox}
\section[le beffroi]{le beffroi — the belfry (town bell tower)}

\label{lesson:FR-C439-beffroi}



\begin{quote}
Before the new one: say the French for the basilica, then the French for the bell tower.
\end{quote}

\subsection*{You'll want to know: le beffroi}
\textbf{le beffroi} — \textquotedblleft{}the belfry (town bell tower)\textquotedblright{}.

\begin{quote}
Du haut du beffroi, la vue sur la ville est superbe.
\end{quote}

\begin{quote}
\emph{From the top of the belfry the view over the town is superb.}
\end{quote}

Many towns in the north of France and in Belgium have a \emph{beffroi}.

\begin{cousinweb}
\emph{beffroi} is usually traced to Middle High German \emph{bercfrit} 'tower that protects'; EtymDB does not list it.
\end{cousinweb}

\begin{grammarlens}[title={What you've built}]
One more word for telling, explaining and describing at B1.
\end{grammarlens}

\subsection*{Your turn}
\begin{itemize}
\raggedright
  \item \emph{Say it:} \emph{le beffroi}
  \item \emph{Say it:} \emph{le beffroi}, once more
  \item \emph{Say it:} a sentence of your own with \emph{le beffroi}
  \item \emph{Recall:} say the French for the basilica, then the French for the bell tower, then say \emph{le beffroi} again
\end{itemize}

\begin{tcolorbox}[breakable,colback=teal!4,colframe=teal!35!black,title={Before you move on}]
What does le beffroi mean? (\textquotedblleft{}The belfry (town bell tower)\textquotedblright{}.) Say it once more.
\end{tcolorbox}
\section[les remparts]{les remparts — the town walls, the ramparts}

\label{lesson:FR-C439-remparts}



\begin{quote}
Before the new one: say the French for the bell tower, then the French for the belfry.
\end{quote}

\subsection*{You'll want to know: les remparts}
\textbf{les remparts} — \textquotedblleft{}the town walls, the ramparts\textquotedblright{}.

\begin{quote}
Le soir, nous faisons le tour des remparts.
\end{quote}

\begin{quote}
\emph{In the evening we walk round the town walls.}
\end{quote}

Mostly plural.

\begin{cousinweb}
\emph{rempart} is built on \emph{remparer} 'to fortify again', from \emph{re-} and \emph{emparer}, which is usually traced to Provençal \emph{amparar} 'to protect', from Latin \emph{ante} and \emph{parāre} (Lewis \& Short). EtymDB does not list it.
\end{cousinweb}

\begin{grammarlens}[title={What you've built}]
One more word for telling, explaining and describing at B1.
\end{grammarlens}

\subsection*{Your turn}
\begin{itemize}
\raggedright
  \item \emph{Say it:} \emph{les remparts}
  \item \emph{Say it:} \emph{les remparts}, once more
  \item \emph{Say it:} a sentence of your own with \emph{les remparts}
  \item \emph{Recall:} say the French for the bell tower, then the French for the belfry, then say \emph{les remparts} again
\end{itemize}

\begin{tcolorbox}[breakable,colback=teal!4,colframe=teal!35!black,title={Before you move on}]
What does les remparts mean? (\textquotedblleft{}The town walls, the ramparts\textquotedblright{}.) Say it once more.
\end{tcolorbox}
\section[(dialogue)]{Retour sur les chapitres 432–439}

\label{lesson:FR-R439-revision-432-439}



\begin{quote}
Say the French for the manager and the inconvenience.
\end{quote}

\subsection*{Your turn}
Say these aloud:

\begin{itemize}
\raggedright
  \item the manager, the inconvenience, to report, the cleaning, faulty
  \item the down payment, the deposit, to compensate, the compensation, to pay damages to
  \item to get one’s bearings, the compass, to spot, the landmark, to mislay
  \item the road sign, the arrow, the crossroads, the marker post, the corner
  \item the shortcut, the detour, the route, to go round, to turn back
  \item the dead end, the alley, the avenue, the boulevard, the passer-by
  \item to stroll, to go for a walk, to walk about, to wander, to pace up and down
  \item the cathedral, the basilica, the bell tower, the belfry, the town walls
\end{itemize}

\begin{tcolorbox}[breakable,colback=teal!4,colframe=teal!35!black,title={Before you move on}]
The manager? (\textbf{le gérant}.) The inconvenience? (\textbf{le désagrément}.) To report? (\textbf{signaler}.) The cleaning? (\textbf{le ménage}.) Faulty? (\textbf{défectueux}.) The down payment? (\textbf{l'acompte}.) The deposit? (\textbf{les arrhes}.) To compensate? (\textbf{dédommager}.) The compensation? (\textbf{le dédommagement}.) To pay damages to? (\textbf{indemniser}.) To get one’s bearings? (\textbf{orienter}.) The compass? (\textbf{la boussole}.) To spot? (\textbf{repérer}.) The landmark? (\textbf{le repère}.) To mislay? (\textbf{égarer}.) The road sign? (\textbf{le panneau}.) The arrow? (\textbf{la flèche}.) The crossroads? (\textbf{le croisement}.) The marker post? (\textbf{la borne}.) The corner? (\textbf{l'angle}.) The shortcut? (\textbf{le raccourci}.) The detour? (\textbf{le détour}.) The route? (\textbf{l'itinéraire}.) To go round? (\textbf{contourner}.) To turn back? (\textbf{rebrousser}.) The dead end? (\textbf{l'impasse}.) The alley? (\textbf{la ruelle}.) The avenue? (\textbf{l'avenue}.) The boulevard? (\textbf{le boulevard}.) The passer-by? (\textbf{le passant}.) To stroll? (\textbf{flâner}.) To go for a walk? (\textbf{balader}.) To walk about? (\textbf{déambuler}.) To wander? (\textbf{errer}.) To pace up and down? (\textbf{arpenter}.) The cathedral? (\textbf{la cathédrale}.) The basilica? (\textbf{la basilique}.) The bell tower? (\textbf{le clocher}.) The belfry? (\textbf{le beffroi}.) The town walls? (\textbf{les remparts}.)
\end{tcolorbox}
